\section{Inclusion proof: construction and verification}
\label{sec:algorithm}

\subsection{Two kinds of revealed value}

Besides the target pair's two raw data elements, every proof step reveals
exactly one 32-byte (in general, $w$-bit) state value per level of
recursion. Before saying what each reveal means, we need two pieces of
notation: the two snapshots of the running hash-chain state taken while a
block is built, and the notion of a \emph{sibling}.

\begin{definition}[Prewrap and ending]
\label{def:prewrap-ending}
Building $\mathrm{block}(a,k)$ (Section~\ref{sec:construction}) is a
prefix of the single, ongoing compression-chain computation for the whole
file, and it ends with one specific compression call: the one that applies
the block's own closing checkpoint re-injection, $\mathrm{MID}(a)$. Write
\[
\mathrm{prewrap}(a,k) = \text{the running state \emph{immediately before}
that final call,}
\]
i.e.\ right after every element and every nested $\mathrm{MID}$ belonging to
$\mathrm{block}(a,k)$ \emph{other than its own closing one} has been hashed,
and write
\[
\mathrm{ending}(a,k) = \C\big(\mathrm{prewrap}(a,k),\ \mathrm{MID}(a)\big)
\]
for the state \emph{immediately after} that call. Both are snapshots of the
same single running state used to build the whole file -- they are not
separately computed side channels -- taken one compression call apart, on
either side of $\mathrm{block}(a,k)$'s own closing step. (A checkpoint
$\CK[a']$ for $a' \ne 0$, Section~\ref{sec:construction}, is always equal
to $\mathrm{ending}(a,k)$ for whichever block happens to finish immediately
before position $a'$ begins; the one exception is $\CK[0] = \mathrm{IV}$,
which is not any block's $\mathrm{ending}$.)
\end{definition}

\paragraph{The generic step.} Before a concrete example, here is the same
pattern once, abstractly, for a single call $\mathrm{block}(a,k)$ with
$k \ge 2$. Every line below is one of exactly two kinds, and the rule is
about the \emph{word}, not the font: a line that says \textbf{emit} writes
bytes to the file and feeds them into the hash chain; every other line
merely \emph{names} a value the running state already has, and writes
nothing. (\emph{Execute} is a third thing, strictly speaking -- a recursive
call that itself expands into a sequence of exactly this same shape, one
level down -- so it is set in plain text, neither bold nor colored.) To
make the two real kinds easy to tell apart at a glance even where the
line is mostly math, \emph{emit} lines are \textbf{bold black}, and every
value-naming line is \snap{shown in gray}:

\begin{enumerate}[itemsep=1pt, parsep=0pt, topsep=4pt]
\item \snap{if not already set: $\CK[a] := $ state}
\item execute $\mathrm{block}(a,\, k{-}1)$
\item execute $\mathrm{block}(a + 2^{k-1},\, k{-}1)$
\item \snap{$\mathrm{prewrap}(a,k) := $ state}
\item \textbf{emit $\mathrm{MID}(a)$}
\item \snap{$\mathrm{ending}(a,k) := $ state}
\end{enumerate}

The base case $k=1$ is the same six steps with two data emissions,
\textbf{emit $D[a]$} and \textbf{emit $D[a{+}1]$}, standing in for the two
recursive calls (steps 2--3). Every call to $\mathrm{block}(\cdot,\cdot)$
made anywhere in the whole construction -- there are exactly $N-2$ of them,
one per $\mathrm{MID}$ emission -- follows one of these two shapes.

The $N=8$ example below is exactly what you get from inlining every one of
the $N-2=6$ calls that construction makes, in the order they actually
execute, down to the individual data elements.

\paragraph{Example: the full build sequence for $N=8$.} To make
Definition~\ref{def:prewrap-ending} concrete, here is the complete,
real sequence of events building the structured file for a small $N=8$
example ($D[0..7]$), exactly as produced by the reference implementation,
using the same \textbf{bold} (emit) / \snap{gray} (value-naming only)
convention as above:

\begin{enumerate}[itemsep=1pt, parsep=0pt, topsep=4pt]
\item \snap{$\CK[0] := \mathrm{IV}$}
\item \textbf{emit $D[0]$}
\item \textbf{emit $D[1]$}
\item \snap{$\mathrm{prewrap}(0,1) :=$ state}
\item \textbf{emit $\mathrm{MID}(0)$} \hfill (bytes $= \mathrm{pad}_b(\CK[0])$)
\item \snap{$\mathrm{ending}(0,1) :=$ state}
\item \snap{$\CK[2] :=$ state} \hfill ($= \mathrm{ending}(0,1)$)
\item \textbf{emit $D[2]$}
\item \textbf{emit $D[3]$}
\item \snap{$\mathrm{prewrap}(2,1) :=$ state}
\item \textbf{emit $\mathrm{MID}(2)$} \hfill (bytes $= \mathrm{pad}_b(\CK[2])$)
\item \snap{$\mathrm{ending}(2,1) :=$ state}
\item \snap{$\mathrm{prewrap}(0,2) :=$ state} \hfill ($= \mathrm{ending}(2,1)$)
\item \textbf{emit $\mathrm{MID}(0)$ (again)} \hfill (bytes $=
  \mathrm{pad}_b(\CK[0])$, same bytes as step 5)
\item \snap{$\mathrm{ending}(0,2) :=$ state}
\item \snap{$\CK[4] :=$ state} \hfill ($= \mathrm{ending}(0,2)$)
\item \textbf{emit $D[4]$}
\item \textbf{emit $D[5]$}
\item \snap{$\mathrm{prewrap}(4,1) :=$ state}
\item \textbf{emit $\mathrm{MID}(4)$} \hfill (bytes $= \mathrm{pad}_b(\CK[4])$)
\item \snap{$\mathrm{ending}(4,1) :=$ state}
\item \snap{$\CK[6] :=$ state} \hfill ($= \mathrm{ending}(4,1)$)
\item \textbf{emit $D[6]$}
\item \textbf{emit $D[7]$}
\item \snap{$\mathrm{prewrap}(6,1) :=$ state}
\item \textbf{emit $\mathrm{MID}(6)$} \hfill (bytes $= \mathrm{pad}_b(\CK[6])$)
\item \snap{$\mathrm{ending}(6,1) :=$ state}
\item \snap{$\mathrm{prewrap}(4,2) :=$ state} \hfill ($= \mathrm{ending}(6,1)$)
\item \textbf{emit $\mathrm{MID}(4)$ (again)} \hfill (bytes $=
  \mathrm{pad}_b(\CK[4])$, same bytes as step 20)
\item \snap{$\mathrm{ending}(4,2) :=$ state} \hfill (this is the file's raw
  chain state, before finalization)
\end{enumerate}

Two things are worth noticing. First, every \snap{gray, non-\textbf{emit}}
line is free: it is a name for a value the running state already has at
that point, not a separate computation and not a separate byte written
anywhere -- Definition~\ref{def:prewrap-ending}'s
$\mathrm{prewrap}$/$\mathrm{ending}$ pair is nothing more than \emph{before}
and \emph{after} labels for step~5 (or step~11, or~20, or~26 -- whichever
$\mathrm{MID}$ call is currently closing). Second, $\mathrm{MID}(0)$ is
emitted twice (steps 5 and
14) and $\mathrm{MID}(4)$ is emitted twice (steps 20 and 29): same
checkpoint bytes both times, but compressed against a different running
state each time, since the same original position can start more than one
nested block (here $a=0$ starts both $\mathrm{block}(0,1)$ and
$\mathrm{block}(0,2)$). This is exactly the mechanism that lets one
checkpoint, computed once, be reused at every level whose block happens to
start at that position -- the same reuse Definition~\ref{def:prewrap-ending}
already alludes to for $\CK[a']$ in general. (14 emitted items total -- 8
data, 6 $\mathrm{MID}$ -- matching $2N-2$ for $N=8$, Section~\ref{sec:construction}.)

\begin{definition}[Sibling]
\label{def:sibling}
At recursion level $k$, block $\mathrm{block}(a,k)$ has two children,
$\mathrm{block}(a, k-1)$ and $\mathrm{block}(a+h, k-1)$ with $h = 2^{k-1}$
(Section~\ref{sec:construction}). The target index $t$ lies in exactly one
of the two; call that one the \emph{target child}. The \textbf{sibling} (at
this level) is the other child -- the one \emph{not} containing $t$.
\end{definition}

Since the sibling does not contain the target, the proof does not recurse
into it: it reveals a single state value that stands in for the sibling's
whole contents. Which value depends on whether the sibling was built
before or after the target child:

\begin{description}
\item[\emph{Left sibling: becomes the new checkpoint.}] The sibling
  $\mathrm{block}(a,k-1)$ is already closed when the target child begins.
  Its closing state $\mathrm{ending}(a,k-1) = \CK_{a+2^{k-1}}$ is revealed
  and used directly as the checkpoint that starts the evaluation of the
  target child.
\item[\emph{Right sibling: combined with the verified value.}] The sibling
  $\mathrm{block}(a+h,k-1)$ is built after the target child. The proof
  reveals its \emph{pre}-closing state $\mathrm{prewrap}(a+h,k-1)$, not its
  ending, and the verifier performs the sibling's closing fold itself,
  using $\mathit{left\_end}$, its own result for the target child, as the
  checkpoint. Revealing $\mathrm{ending}(a+h,k-1)$ instead would make the
  final check independent of $\mathit{left\_end}$, so any target data below
  this level would verify (Remark~\ref{rem:prewrap-bug}). For $k-1 \ge 2$,
  $\mathrm{prewrap}(a+h,k-1)$ is the ending of the sibling's own right child
  and carries the label $\CK_{[a+h+2^{k-2},\,a+h+2^{k-1}-1]}$
  (Section~\ref{sec:construction}); for $k-1=1$ it is unlabelled.
\end{description}

Concretely, the fold-in step in the second case computes
\begin{equation}
\C\big(\mathrm{prewrap}(a{+}h,k{-}1),\ \mathrm{pad}_b(\mathit{left\_end})\big),
\label{eq:foldin}
\end{equation}
and in a genuine proof $\mathit{left\_end} = \mathrm{ending}(a,k-1) =
\CK[a+h]$, the checkpoint that $\mathrm{block}(a+h,k-1)$'s closing
$\mathrm{MID}(a+h)$ re-injects, so by Definition~\ref{def:prewrap-ending}
\begin{equation}
\C\big(\mathrm{prewrap}(a{+}h,k{-}1),\ \mathrm{pad}_b(\CK[a+h])\big)
= \mathrm{ending}(a{+}h,k{-}1)
\end{equation}
exactly: the verifier recomputes the sibling's closing value from a
checkpoint it has itself just verified.

\paragraph{Closing $\mathrm{block}(a,k)$ itself.} After both children have
been accounted for, in left-to-right order, the state is
$v = \mathrm{prewrap}(a,k)$ in either case, and one further call closes
$\mathrm{block}(a,k)$ with its own checkpoint $\CK[a]$, passed down from
the level above:
\begin{equation}
\C\big(v,\ \mathrm{pad}_b(\CK[a])\big)
= \C\big(\mathrm{prewrap}(a,k),\ \mathrm{MID}(a)\big)
= \mathrm{ending}(a,k).
\label{eq:own-close}
\end{equation}
This is the last line of \textsc{EvalBlock} below,
\texttt{return C(right\_end, pad\_b(ckpt))}. Every level performs
Equation~\eqref{eq:own-close} once, in addition to resolving the sibling.
The two folds are easy to conflate in the pseudocode, but
Equation~\eqref{eq:foldin} closes the \emph{sibling} with $\CK[a{+}h]$,
while Equation~\eqref{eq:own-close} closes $\mathrm{block}(a,k)$ with
$\CK[a]$, one level up.

\subsection{Building a proof}

\begin{algorithm}[H]
\caption{$\textsc{BuildProof}(a_0, L)$ -- prover, given target pair start $a_0$}
\begin{lstlisting}[style=pseudo]
function Recurse(a, k, t):      # t = target index, in [a, a+2^k)
    if k = 1: return             # base case: no reveal here,
                                  # base pair always sent raw
    h := 2^(k-1)
    if t < a + h:                # target in left child
        Recurse(a, k-1, t)
        reveal prewrap(a+h, k-1)  # right sibling, not yet closed
    else:                         # target in right child
        reveal ending(a, k-1)     # left sibling, already closed
        Recurse(a+h, k-1, t)

h_top := 2^(L-1)
if a0 < h_top:
    Recurse(0, L-1, a0)
    reveal prewrap(h_top, L-1)
else:
    reveal ending(0, L-1)
    Recurse(h_top, L-1, a0)

return (D[a0], D[a0+1], reveals in the order emitted above)
\end{lstlisting}
\end{algorithm}

The proof carries $a_0$'s two data elements, not $a_0$ itself: the
verifier (Algorithm~3.3 below) always has its own trusted target
$\mathit{index}$ and derives the pair start from \emph{that}, never from
anything the proof claims -- $a_0$ is prover-side bookkeeping used to
build the reveal sequence, not part of the wire format.

\begin{lemma}[Index-independent proof size, for fixed $N$]
For every index $0 \le a < N = 2^L$, this procedure emits exactly $L-1$
reveals -- the same count for every index at that $N$, not merely an
average over indices. (Proof size is, as usual, still $O(\log N)$ in $N$
itself; the point here is that it does not additionally vary with
\emph{which} index is proven once $N$ is fixed.)
\end{lemma}

This follows directly from the algorithm: each of the $L-1$ recursive levels
above the base case emits exactly one reveal, regardless of which branch is
taken. (Section~\ref{sec:algorithm-cost} discusses \emph{verifier compute
cost}, which is \emph{not} constant across indices, even though proof size
is.)

\paragraph{Prover time and storage.} \textsc{BuildProof} as written treats
$\mathrm{ending}(a,k)$ and $\mathrm{prewrap}(a,k)$ as values simply
available to read off -- but, unlike the checkpoints $\CK[a]$, neither is a
literal byte range stored anywhere in the structured file
(Definition~\ref{def:prewrap-ending}): both are internal running-chain-
state snapshots that existed only transiently while the file was originally
built. Section~\ref{sec:construction}'s $O(\log N)$ working-memory claim is
for the \emph{build} step alone, and explicitly discards each checkpoint
once its block closes; it says nothing about the cost of later answering an
arbitrary proof request. Two honest options exist, with a real space/time
trade-off between them, and a deployment needs to pick one explicitly:
\begin{itemize}
\item \textbf{Cache every snapshot at build time.} Retain all $O(N)$
  intermediate $\mathrm{ending}$/$\mathrm{prewrap}$ values (one 32-byte, or
  in general $w$-bit, entry per $\mathrm{MID}$ emission) alongside the
  file, rather than discarding them. \textsc{BuildProof} then costs
  $O(\log N)$ time per request, matching proof \emph{size}, at the cost of
  $O(N)$ storage held indefinitely by whoever answers proof requests (in
  addition to the file itself). This is what the reference implementation
  does, for simplicity, exactly as Section~\ref{sec:construction} already
  notes for the analogous build-time choice.
\item \textbf{Recompute on demand.} Keep only the file bytes and answer
  each proof request by re-running the hash chain from the nearest
  available checkpoint up to the needed snapshot. Without any cached
  snapshots this costs up to $O(N)$ compression calls per proof request --
  no better than a full re-hash -- though a partial cache (e.g.\ one
  snapshot every $O(\sqrt{N})$ positions) trades a smaller, tunable storage
  budget against a correspondingly higher, but still sublinear, per-request
  recomputation cost.
\end{itemize}
Either way, the $O(\log N)$ figures emphasized elsewhere in this paper are
proof \emph{size} and verifier \emph{compute}, not prover storage: a
FlatMerkle prover that wants every request answered in $O(\log N)$ time,
matching the verifier's own cost, needs $O(N)$ storage to do it, the same
asymptotic footprint an equivalent Merkle-tree prover needs to keep the
whole tree (not just the root) on hand.

\subsection{Verifying a proof}

\begin{algorithm}[H]
\caption{$\textsc{EvalBlock}(a, k, t, \mathit{ckpt})$ -- verifier, symmetric to \textsc{Recurse}}
\label{alg:evalblock}
\begin{lstlisting}[style=pseudo]
function EvalBlock(a, k, t, ckpt):
    if k = 1:
        s := C(C(ckpt, D0), D1)      # the proof's base pair
        return C(s, pad_b(ckpt))     # re-inject checkpoint,
                                      # closing the base pair
    h := 2^(k-1)
    if t < a + h:                     # left branch
        left_end  := EvalBlock(a, k-1, t, ckpt)
        sib_prewrap := read next reveal
        right_end := C(sib_prewrap, pad_b(left_end))
                                       # fold in the value just verified
        return C(right_end, pad_b(ckpt))
    else:                              # right branch
        left_end  := read next reveal
        right_end := EvalBlock(a+h, k-1, t, left_end)
        return C(right_end, pad_b(ckpt))
\end{lstlisting}
\end{algorithm}

\begin{algorithm}[H]
\caption{$\textsc{Verify}(N, \Delta, \mathit{index}, \mathit{data}, \pi)$ -- top-level verifier}
\begin{lstlisting}[style=pseudo]
function Verify(N, Delta, index, data, pi):
    if N < 4 or N is not a power of two: return false   # domain: N = 2^L, L >= 2
    if index < 0 or index >= N: return false   # index out of range
    (D0, D1, reveals) := pi                    # reveals: a shared queue,
                                                # consumed front-to-back by
                                                # every "read next reveal"
                                                # line below and inside
                                                # EvalBlock (Section 3.3)
    a0 := index - (index mod 2)                # pair start, derived from
                                                # the verifier's OWN index --
                                                # NEVER taken from pi
    if index = a0: expected := D0 else: expected := D1
    if data != expected: return false           # supplied data disagrees with proof
    L := log2(N)
    h_top := 2^(L-1)
    if a0 < h_top:
        left_end := EvalBlock(0, L-1, a0, IV)     # consumes from reveals
        if reveals is empty: return false          # ran out before the top fold
        sib_prewrap := read next reveal
        raw_state := C(sib_prewrap, pad_b(left_end))
    else:
        if reveals is empty: return false
        left_end := read next reveal
        raw_state := EvalBlock(h_top, L-1, a0, left_end)
    if reveals is not empty: return false        # leftover reveals: malformed
    return finalize(raw_state, N) = Delta          # public length-padding + compare
\end{lstlisting}
\end{algorithm}

\begin{remark}[A natural-looking but insecure alternative: trusting a
proof-supplied pair start]
\label{rem:a0-bug}
A verifier might instead accept a pair start $a_0$ as part of the proof
itself, $\pi = (a_0, D_0, D_1, \mathit{reveals})$, and branch on whether
$\mathit{index} = a_0$ or $\mathit{index} = a_0+1$ to decide which of
$D_0, D_1$ is ``the'' claimed data -- exactly as \textsc{BuildProof}
computes $a_0$ internally. This is a genuine soundness hole, and a
collision-free one: \textsc{EvalBlock}'s recursion branches only on
whether the recursion target lies below a block's midpoint, and every
branch point above the base case compares against a power-of-two
threshold $h = 2^{k-1} \ge 2$, so the least-significant bit of the
recursion target never affects any branch decision -- indices $i-1$ and
$i$ (for odd $i$) take an \emph{identical} recursion path. An honest
proof for index $i-1$ carries $D_0 = D[i-1]$, and a verifier that trusts a
proof-supplied $a_0$ can be asked to check $(\mathit{index}=i,\
\mathit{data}=D[i-1], \pi)$ with $a_0$ relabelled to $i$ instead of $i-1$
(same $D_0, D_1$, same reveals, unchanged): the trusting verifier sees
$\mathit{index}=a_0$, sets $\mathit{expected} := D_0 = D[i-1]$, and
accepts -- reaching the genuine digest with no free-start collision at
all, since the recursion itself never used $a_0$ for anything but this one
comparison. The construction above avoids this by design: $a_0$ is never
read from $\pi$, only derived from the verifier's own $\mathit{index}$,
which the verifier is asking about in the first place and has no reason
to distrust. (The reference implementations already follow this
discipline -- \texttt{verify\_proof} derives its pair start solely from
the caller-supplied target index, and the wire-level \texttt{Proof} object
carries no separate pair-start field at all.)
\end{remark}

\textsc{EvalBlock}'s ``read next reveal'' lines share the same queue and
must reject if it is empty: too few reveals is as malformed as too many,
and an implementation that substitutes a default value for a missing
reveal reopens a gap of the same kind as Remark~\ref{rem:prewrap-bug}'s.
\textsc{Verify} first enforces the domain $N = 2^L$, $L \ge 2$
(Section~\ref{sec:construction}; for $N=2$ there is no top-level split),
then $\mathit{index} \in [0,N)$, derives $a_0$ from $\mathit{index}$
alone, checks $\mathit{data}$ against the element of the base pair that
$\mathit{index}$ selects, mirrors \textsc{BuildProof}'s top-level split
(there is no $\mathrm{MID}$ joining the two root blocks; the finalization
plays that role), rejects missing or leftover reveals, and only then
compares the finalized chain state with the digest. Both reference
implementations perform all of these checks; a re-implementation of
\textsc{EvalBlock}'s recursive core alone should add them.

\begin{remark}[A natural-looking but insecure alternative]
\label{rem:prewrap-bug}
A verifier might be tempted to compute the left branch's
$\mathit{left\_end}$ recursively (as required) but then discard it,
returning
\[
\C(\mathit{right\_end}, \mathrm{pad}_b(\mathit{ckpt}))
\]
without ever folding $\mathit{left\_end}$ into the check. This is a
genuine soundness hole: tampering with any value inside the left recursion
goes undetected, because the final check never actually depends on it.
The construction above avoids this by design (left branch of Algorithm~\ref{alg:evalblock}):
\[
\mathit{right\_end} := \C(\mathit{sib\_prewrap}, \mathrm{pad}_b(\mathit{left\_end})).
\]
The verifier must combine
the just-verified $\mathit{left\_end}$ with the revealed sibling
\emph{pre-close} state, not simply trust a separately revealed
\emph{post-close} state for the sibling and ignore what it just computed.
We recommend testing for exactly this class of omission directly --
flipping a single bit in every individually revealed value and confirming
rejection -- since it is the kind of gap that is easy to miss by
inspection alone.
\end{remark}

\subsection{Verifier compute cost varies by index; proof size does not}
\label{sec:algorithm-cost}

Although every proof has exactly $L-1$ reveals, the number of
compression-function calls the \emph{verifier} performs to check a proof is
not constant: it depends on how often the target's recursion path takes the
left branch (which costs 2 calls at that level: the fold, then the
checkpoint re-injection) versus the right branch (which costs 1 call: only
the checkpoint re-injection, since the sibling's closed value is used
as-is). For $N=32$ ($L=5$), measured directly against the real verifier
implementation:

\begin{table}[htbp]
\centering
\begin{tabular}{@{}lccccc@{}}
\toprule
Index & 0 & 8 & 15 & 23 & 30 \\
\midrule
Compression calls (incl.\ finalization) & 11 & 10 & 8 & 8 & 7 \\
\bottomrule
\end{tabular}
\caption{Verifier compression-function calls by target index, $N=32$. The
minimum, 7, is achieved by indices that take the cheap (right) branch at
every level; the maximum, 11, by index 0, which takes the left branch at
every level. In general the range is $[L+2,\ 2L+1]$ (including one constant
finalization call): the ratio $(2L+1)/(L+2)$ approaches, but never reaches,
$2\times$ as $L \to \infty$, and is strictly less than $2\times$ for every
finite $L$ (here, $L=5$: $11/7 \approx 1.57\times$) -- bounded, not
asymptotically growing relative to $L$.}
\label{tab:verifier-cost}
\end{table}

This asymmetry is a lever, not just a cost -- but the right way to describe
it is by \emph{bit pattern}, not by index magnitude. The exact cost is
$L + 2 + \#\{\text{zero bits among bits } 1..L{-}1 \text{ of the target}\}$,
which depends only on which of those bits are set, not on how large the
index is: for $N=32$, index 16 ($\texttt{10000}_2$) costs 10 -- almost as
expensive as index 0's 11, and no cheaper than index 8's 10 -- despite
being numerically closer to the ``high-index end'' than either. Index 30
($\texttt{11110}_2$) gets the minimum cost of 7 not because it is large,
but because its high-order bits happen to all be 1. The correct lever is
therefore: if the elements most likely to be queried are known in advance,
placing them at indices whose high-order bits are all 1 -- concretely, the
last few elements of the file -- favors the cheap, right-branch-heavy
verification path and reduces their average verification cost for free,
with no change to the construction or to the constant proof-size
guarantee. Since the whole range is bounded, $[L+2,\ 2L+1]$, the payoff of
this placement is at most a factor of roughly 2 and does not grow with
$N$ -- a genuine but modest optimization, not an asymptotic one.

\paragraph{A proof always discloses its whole base pair.}
\label{par:pair-disclosure}
\textsc{BuildProof} returns $D[a_0]$ and $D[a_0{+}1]$ unconditionally, where
$a_0 = a - (a \bmod 2)$ -- so a proof requested for target index $a$ always
discloses \emph{both} raw elements of its pair, not just $D[a]$: proving
index 7 necessarily also reveals element 6, whether or not anyone asked
about element 6. This is a real, semantic difference from an ordinary
Merkle proof, whose leaf hash need not expose any neighboring leaf, and it
is worth stating plainly rather than leaving implicit. It also affects the
size comparison in Section~\ref{sec:examples}'s table and
Section~\ref{sec:optimal-size}: with element size equal to the state width,
$D=W$, FlatMerkle's proof size $2D + (L-1)W$ and a same-shape Merkle
proof's size $LW + D$ come out numerically \emph{equal} -- but FlatMerkle's
byte count buys one fewer sibling hash \emph{and} one extra disclosed
element, not a strictly smaller disclosure. Anyone relying on FlatMerkle
proofs where the neighboring element's privacy matters should account for
this explicitly.
